% SPDX-FileCopyrightText: 2026 Vitor Mattos
% SPDX-License-Identifier: CC-BY-SA-4.0
\documentclass[10pt,aspectratio=169,t]{beamer}

\usepackage[utf8]{inputenc}
\usepackage[T1]{fontenc}
\usepackage[brazilian]{babel}
\usepackage{lmodern}
\usepackage{microtype}
\usepackage{xcolor}
\usepackage{tikz}
\usetikzlibrary{arrows.meta,positioning,calc,fit}
\usepackage{hyperref}

% Identidade visual inspirada em apresentações acadêmicas da UTFPR.
% Não usa marca ou ativo institucional oficial.
\definecolor{utfpryellow}{HTML}{F4C300}
\definecolor{utfprdark}{HTML}{1D1D1F}
\definecolor{utfprblue}{HTML}{1F5A94}
\definecolor{utfprteal}{HTML}{277C78}
\definecolor{muted}{HTML}{6B7280}
\definecolor{softgray}{HTML}{F3F4F6}
\definecolor{softblue}{HTML}{E8F0F7}
\definecolor{softyellow}{HTML}{FFF7D6}
\definecolor{softgreen}{HTML}{E7F1EA}
\definecolor{softred}{HTML}{F8E6E6}

\setbeamertemplate{navigation symbols}{}
\setbeamercolor{normal text}{fg=utfprdark,bg=white}
\setbeamercolor{frametitle}{fg=utfprdark,bg=white}
\setbeamerfont{frametitle}{series=\bfseries,size=\Large}
\setbeamerfont{title}{series=\bfseries,size=\huge}
\setbeamertemplate{itemize item}{\color{utfpryellow}$\blacktriangleright$}
\setbeamertemplate{itemize subitem}{\color{utfprblue}$\bullet$}

\setbeamertemplate{frametitle}{%
  \nointerlineskip
  \begin{beamercolorbox}[wd=\paperwidth,ht=0.74cm,dp=0.18cm,leftskip=0.55cm,rightskip=0.5cm]{frametitle}
    \insertframetitle
  \end{beamercolorbox}
  \vspace{-0.10cm}
  \hspace{0.55cm}\textcolor{utfpryellow}{\rule{2.6cm}{2.4pt}}
}

\setbeamertemplate{footline}{%
  \leavevmode
  \hbox{%
  \begin{beamercolorbox}[wd=.80\paperwidth,ht=2.4ex,dp=.9ex,leftskip=.55cm]{author in head/foot}
    \scriptsize UTFPR -- Software Livre \hspace{0.7em}|\hspace{0.7em} Vitor Mattos
  \end{beamercolorbox}%
  \begin{beamercolorbox}[wd=.20\paperwidth,ht=2.4ex,dp=.9ex,rightskip=.55cm plus1fil]{date in head/foot}
    \scriptsize \insertframenumber/\inserttotalframenumber
  \end{beamercolorbox}}
}

\newcommand{\source}[1]{\vfill{\scriptsize\color{muted}#1}}
\newcommand{\bigclaim}[1]{\begin{center}\Large\bfseries #1\end{center}}
\newcommand{\pill}[2]{\tikz[baseline=(n.base)]\node[rounded corners=4pt,fill=#1,inner xsep=6pt,inner ysep=3pt](n){\scriptsize #2};}
\newcommand{\tagbox}[3]{\node[rounded corners=7pt,fill=#1,draw=#2,thick,inner sep=7pt,text width=2.7cm,align=center] {#3};}
\newcommand{\stage}[4]{\node[rounded corners=8pt,fill=#2,draw=#3,thick,minimum width=3.1cm,minimum height=1.1cm,align=center] (#1) {#4};}

\title{Do software proprietário ao ecossistema aberto}
\subtitle{Leitura crítica de Kilamo et al. (2012)}
\author{Vitor Mattos}
\institute{Disciplina de Software Livre -- UTFPR}
\date{2026}

\begin{document}

\begin{frame}[plain]
  \vspace{0.35cm}
  {\huge\bfseries Do software proprietário ao ecossistema aberto}\par
  \vspace{0.15cm}
  {\large Leitura crítica de Kilamo et al. (2012)}\par
  \vspace{0.65cm}
  \begin{center}
  \begin{tikzpicture}[node distance=0.72cm]
    \stage{a}{utfprdark}{utfprdark}{\color{white}produto\\fechado}
    \stage{b}{softyellow}{utfpryellow}{abertura}
    \stage{c}{softgreen}{utfprteal}{ecossistema\\aberto}
    \draw[-{Stealth[length=3mm]},line width=1.1pt,utfpryellow] (a)--(b);
    \draw[-{Stealth[length=3mm]},line width=1.1pt,utfprteal] (b)--(c);
  \end{tikzpicture}
  \end{center}
  \vspace{0.45cm}
  {\small Terhi Kilamo, Imed Hammouda, Tommi Mikkonen, Timo Aaltonen. \textit{From proprietary to open source: Growing an open source ecosystem}. \textit{Journal of Systems and Software}, 85(7), 1467--1478, 2012.}
  \source{DOI: \href{https://doi.org/10.1016/j.jss.2011.06.071}{10.1016/j.jss.2011.06.071} \hspace{0.5em}|\hspace{0.5em} Fonte: \href{https://www.sciencedirect.com/science/article/pii/S0164121211001683}{ScienceDirect}}
\end{frame}

\begin{frame}{Abrir o código não basta}
  \vspace{0.55cm}
  \bigclaim{O problema do artigo é transformar abertura em continuidade.}
  \vspace{0.45cm}
  \begin{center}
  \begin{tikzpicture}[node distance=0.48cm]
    \stage{code}{softgray}{muted}{código}
    \stage{entry}{softblue}{utfprblue}{entrada\\de pessoas}
    \stage{rules}{softyellow}{utfpryellow}{coordenação}
    \stage{eco}{softgreen}{utfprteal}{ecossistema}
    \draw[-{Stealth[length=2.5mm]},line width=1pt] (code)--(entry);
    \draw[-{Stealth[length=2.5mm]},line width=1pt] (entry)--(rules);
    \draw[-{Stealth[length=2.5mm]},line width=1pt] (rules)--(eco);
  \end{tikzpicture}
  \end{center}
  \vspace{0.35cm}
  \begin{center}\small\color{muted}Gatilho de fala: repositório público não equivale a comunidade sustentável.\end{center}
  \source{O artigo parte da lacuna de diretrizes para abrir software proprietário e sustentar uma comunidade FLOSS.}
\end{frame}

\begin{frame}{Abertura é um problema de várias dimensões}
  \vspace{0.45cm}
  \begin{center}
  \begin{tikzpicture}[node distance=0.34cm]
    \tagbox{s}{softblue}{utfprblue}{software}
    \tagbox{i}{softgray}{muted}{infraestrutura}
    \tagbox{p}{softyellow}{utfpryellow}{processo}
    \tagbox{l}{softred}{muted}{legalidade}
    \tagbox{m}{softgreen}{utfprteal}{marketing}
    \tagbox{c}{softblue}{utfprblue}{comunidade}
    \node[fit=(s)(i)(p)(l)(m)(c),draw=utfprdark,rounded corners=10pt,inner sep=10pt,label={[font=\small\bfseries]above:seis frentes da abertura}] {};
  \end{tikzpicture}
  \end{center}
  \vspace{0.55cm}
  \begin{center}\Large\bfseries A contribuição do artigo é organizar essas frentes em um processo.\end{center}
  \source{Software, infraestrutura, processo, legalidade, marketing e comunidade aparecem como dimensões da abertura.}
\end{frame}

\begin{frame}{Comunidade tem camadas}
  \vspace{0.25cm}
  \begin{center}
  \begin{tikzpicture}[scale=0.88]
    \fill[softgreen] (0,0) circle (2.6);
    \fill[softyellow] (0,0) circle (2.05);
    \fill[softblue] (0,0) circle (1.5);
    \fill[softred] (0,0) circle (0.95);
    \fill[utfprdark] (0,0) circle (0.42);
    \node[text=white,font=\bfseries\scriptsize] at (0,0) {líder};
    \node[font=\bfseries\scriptsize] at (0,0.82) {núcleo};
    \node[font=\bfseries\scriptsize] at (0,1.35) {ativos};
    \node[font=\bfseries\scriptsize] at (0,1.9) {periféricos};
    \node[font=\bfseries\scriptsize] at (0,2.45) {usuários};
    \node[align=left,text width=4.4cm] at (5.0,0.6) {\Large\bfseries A empresa precisa formar o centro antes de esperar a periferia.};
  \end{tikzpicture}
  \end{center}
  \source{Modelo em cebola de comunidades FLOSS, retomado pelos autores para discutir papéis e engajamento.}
\end{frame}

\begin{frame}{O artigo pergunta como sistematizar a abertura}
  \vspace{0.35cm}
  \begin{center}
  \begin{tikzpicture}[node distance=0.35cm]
    \node[rounded corners=7pt,fill=softblue,draw=utfprblue,thick,text width=2.1cm,align=center,minimum height=1.0cm] (q1) {critérios\\\scriptsize prontidão};
    \node[rounded corners=7pt,fill=softyellow,draw=utfpryellow,thick,right=of q1,text width=2.1cm,align=center,minimum height=1.0cm] (q2) {atores\\\scriptsize plano};
    \node[rounded corners=7pt,fill=softgray,draw=muted,thick,right=of q2,text width=2.1cm,align=center,minimum height=1.0cm] (q3) {dados\\\scriptsize avaliação};
    \node[rounded corners=7pt,fill=softgreen,draw=utfprteal,thick,right=of q3,text width=2.1cm,align=center,minimum height=1.0cm] (q4) {uso\\\scriptsize resultados};
    \node[rounded corners=7pt,fill=softred,draw=muted,thick,right=of q4,text width=2.1cm,align=center,minimum height=1.0cm] (q5) {medição\\\scriptsize comunidade};
  \end{tikzpicture}
  \end{center}
  \vspace{0.75cm}
  \begin{center}\Large\bfseries A pergunta não é ideológica; é operacional.\end{center}
  \source{Os autores listam cinco questões: avaliar prontidão, planejar, coletar dados, usar resultados e medir a comunidade nascente.}
\end{frame}

\begin{frame}{Método: proximidade com a prática}
  \vspace{0.35cm}
  \begin{center}
  \begin{tikzpicture}[node distance=0.55cm]
    \node[rounded corners=7pt,fill=softgray,draw=muted,thick,minimum width=2.9cm,minimum height=0.92cm] (lit) {literatura};
    \node[rounded corners=7pt,fill=softgray,draw=muted,thick,right=of lit,minimum width=2.9cm,minimum height=0.92cm] (work) {workshops};
    \node[rounded corners=7pt,fill=softgray,draw=muted,thick,right=of work,minimum width=2.9cm,minimum height=0.92cm] (cases) {4 empresas};
    \node[rounded corners=8pt,fill=softyellow,draw=utfpryellow,thick,below=0.8cm of work,minimum width=4.2cm,minimum height=1.0cm] (model) {OSCOMM};
    \draw[-{Stealth[length=2.4mm]},line width=0.9pt] (lit)--(model);
    \draw[-{Stealth[length=2.4mm]},line width=0.9pt] (work)--(model);
    \draw[-{Stealth[length=2.4mm]},line width=0.9pt] (cases)--(model);
  \end{tikzpicture}
  \end{center}
  \vspace{0.55cm}
  \begin{columns}[T,onlytextwidth]
    \column{0.48\textwidth}\centering\pill{softgreen}{força}\par\small contexto industrial real
    \column{0.48\textwidth}\centering\pill{softred}{risco}\par\small evidência pouco controlada
  \end{columns}
  \source{Abordagem construtiva e pesquisa-ação com Nokia, Sesca, IT Mill e Gurux.}
\end{frame}

\begin{frame}{OSCOMM: três momentos da abertura}
  \vspace{0.55cm}
  \begin{center}
  \begin{tikzpicture}[node distance=0.75cm]
    \node[rounded corners=8pt,fill=softblue,draw=utfprblue,thick,text width=3.15cm,align=center,minimum height=1.38cm] (r3) {\Large R3\\\small prontidão};
    \node[rounded corners=8pt,fill=softyellow,draw=utfpryellow,thick,right=of r3,text width=3.15cm,align=center,minimum height=1.38cm] (eng) {\Large engenharia\\\small abertura};
    \node[rounded corners=8pt,fill=softgreen,draw=utfprteal,thick,right=of eng,text width=3.15cm,align=center,minimum height=1.38cm] (bulb) {\Large BULB\\\small acompanhamento};
    \draw[-{Stealth[length=3mm]},line width=1.1pt] (r3)--(eng);
    \draw[-{Stealth[length=3mm]},line width=1.1pt] (eng)--(bulb);
  \end{tikzpicture}
  \end{center}
  \vspace{0.55cm}
  \bigclaim{Antes, durante e depois da liberação.}
  \source{O OSCOMM organiza avaliação de prontidão, engenharia de abertura e acompanhamento da comunidade.}
\end{frame}

\begin{frame}{R3: antes de abrir, diagnosticar}
  \vspace{0.45cm}
  \begin{center}
  \begin{tikzpicture}[node distance=0.42cm]
    \node[rounded corners=8pt,fill=softblue,draw=utfprblue,thick,text width=2.55cm,align=center,minimum height=1.0cm] (sw) {software\\\scriptsize código, arquitetura, qualidade};
    \node[rounded corners=8pt,fill=softgreen,draw=utfprteal,thick,right=of sw,text width=2.55cm,align=center,minimum height=1.0cm] (co) {comunidade\\\scriptsize propósito, usuários, parceiros};
    \node[rounded corners=8pt,fill=softred,draw=muted,thick,right=of co,text width=2.55cm,align=center,minimum height=1.0cm] (le) {legalidade\\\scriptsize copyright, licença, marca};
    \node[rounded corners=8pt,fill=softyellow,draw=utfpryellow,thick,right=of le,text width=2.55cm,align=center,minimum height=1.0cm] (au) {autoridade\\\scriptsize cultura, processo, suporte};
  \end{tikzpicture}
  \end{center}
  \vspace{0.55cm}
  \begin{center}\Large\bfseries O R3 não decide sozinho; ele revela gargalos.\end{center}
  \source{As quatro dimensões do R3 têm peso igual, mas os itens internos recebem pesos diferentes.}
\end{frame}

\begin{frame}{Engenharia para abertura: transformar diagnóstico em trabalho}
  \vspace{0.30cm}
  \begin{center}
  \begin{tikzpicture}[node distance=0.28cm]
    \node[rounded corners=7pt,fill=softgray,draw=muted,thick,text width=2.35cm,align=center,minimum height=0.9cm] (ref) {refatorar};
    \node[rounded corners=7pt,fill=softblue,draw=utfprblue,thick,right=of ref,text width=2.35cm,align=center,minimum height=0.9cm] (doc) {documentar};
    \node[rounded corners=7pt,fill=softyellow,draw=utfpryellow,thick,right=of doc,text width=2.35cm,align=center,minimum height=0.9cm] (infra) {infraestrutura};
    \node[rounded corners=7pt,fill=softgreen,draw=utfprteal,thick,right=of infra,text width=2.35cm,align=center,minimum height=0.9cm] (proc) {processo};
    \node[rounded corners=7pt,fill=softred,draw=muted,thick,below=0.42cm of doc,text width=2.35cm,align=center,minimum height=0.9cm] (lic) {licença};
    \node[rounded corners=7pt,fill=softblue,draw=utfprblue,thick,right=of lic,text width=2.35cm,align=center,minimum height=0.9cm] (comm) {entrada};
  \end{tikzpicture}
  \end{center}
  \vspace{0.55cm}
  \bigclaim{A abertura consome esforço antes de produzir comunidade.}
  \source{A fase de engenharia trabalha os problemas encontrados na avaliação de prontidão.}
\end{frame}

\begin{frame}{Acompanhamento: medir se a comunidade existe}
  \vspace{0.35cm}
  \begin{columns}[T,onlytextwidth]
    \column{0.50\textwidth}
      \begin{tikzpicture}[node distance=0.25cm]
        \node[rounded corners=6pt,fill=softblue,draw=utfprblue,thick,text width=4.6cm,align=center] {downloads};
        \node[rounded corners=6pt,fill=softgray,draw=muted,thick,text width=4.6cm,align=center,below=of current bounding box.south] {bugs e pedidos};
      \end{tikzpicture}
    \column{0.50\textwidth}
      \begin{tikzpicture}[node distance=0.25cm]
        \node[rounded corners=6pt,fill=softyellow,draw=utfpryellow,thick,text width=4.6cm,align=center] {comunicação};
        \node[rounded corners=6pt,fill=softgreen,draw=utfprteal,thick,text width=4.6cm,align=center,below=of current bounding box.south] {contribuições};
      \end{tikzpicture}
  \end{columns}
  \vspace{0.70cm}
  \begin{center}\Large\bfseries A comunidade vira objeto de observação contínua.\end{center}
  \source{O modelo BULB coleta sinais como atividade, tamanho da comunidade, comunicação, downloads e contribuições.}
\end{frame}

\begin{frame}{Gurux: o custo concreto da abertura}
  \vspace{0.35cm}
  \begin{columns}[T,onlytextwidth]
    \column{0.25\textwidth}\centering\pill{softgray}{avaliação}\par\Huge 4\par\small pessoas
    \column{0.25\textwidth}\centering\pill{softyellow}{preparação}\par\Huge 6\par\small meses
    \column{0.25\textwidth}\centering\pill{softblue}{equipe}\par\Huge 6\par\small tempo integral
    \column{0.25\textwidth}\centering\pill{softgreen}{comunidade}\par\Huge 200\par\small usuários
  \end{columns}
  \vspace{0.65cm}
  \begin{center}\Large\bfseries Documentação e reescrita custaram quase o dobro do previsto.\end{center}
  \source{Caso Gurux: abertura em 2009, preparação de infraestrutura, documentação, exemplos, fórum e acompanhamento.}
\end{frame}

\begin{frame}{O artigo é forte como guia, fraco como prova}
  \vspace{0.40cm}
  \begin{columns}[T,onlytextwidth]
    \column{0.48\textwidth}
      \begin{tikzpicture}\node[rounded corners=8pt,fill=softgreen,draw=utfprteal,thick,inner sep=10pt,text width=5.0cm]{\bfseries Forças\par\vspace{0.2cm}\small visão integrada\par utilidade industrial\par modelo operacional\par linguagem para decisão};\end{tikzpicture}
    \column{0.48\textwidth}
      \begin{tikzpicture}\node[rounded corners=8pt,fill=softred,draw=muted,thick,inner sep=10pt,text width=5.0cm]{\bfseries Limites\par\vspace{0.2cm}\small poucos casos\par dados confidenciais\par viés de participação\par causalidade fraca};\end{tikzpicture}
  \end{columns}
  \vspace{0.55cm}
  \begin{center}\large\bfseries Ele organiza o problema melhor do que comprova a solução.\end{center}
  \source{Os autores reconhecem que há dados reais, mas ainda insuficientes para evidência conclusiva.}
\end{frame}

\begin{frame}{Implicação prática: construir comunidade é engenharia}
  \vspace{0.45cm}
  \begin{center}
  \begin{tikzpicture}[node distance=0.47cm]
    \stage{c1}{utfprdark}{utfprdark}{\color{white}código}
    \stage{c2}{softblue}{utfprblue}{qualidade}
    \stage{c3}{softyellow}{utfpryellow}{entrada}
    \stage{c4}{softgreen}{utfprteal}{governança}
    \draw[-{Stealth[length=2.5mm]},line width=1pt] (c1)--(c2);
    \draw[-{Stealth[length=2.5mm]},line width=1pt] (c2)--(c3);
    \draw[-{Stealth[length=2.5mm]},line width=1pt] (c3)--(c4);
  \end{tikzpicture}
  \end{center}
  \vspace{0.60cm}
  \bigclaim{A lição para projetos como LibreSign: abertura precisa virar processo social sustentável.}
  \source{A experiência prática funciona aqui como lente crítica, não como evidência adicional do estudo.}
\end{frame}

\end{document}
